% =====================================================================
% Section 5: Experiments
% =====================================================================
%
% Status: Draft Section 5 Experiments v2 (2026-05-17)
% Revision note (v2 vs v1):
%   - Per advisor: borrow HGRAG/Relink experiments style.
%   - Baselines now grouped into 5 categories (No Retrieval /
%     Sparse--Dense / Hierarchical Tree / Graph--Hypergraph /
%     Path--Chain Graph) and the
%     main table prepends a Category column (HGRAG/Relink style).
%   - Subsections renamed and rewritten as RQ1--RQ3 (Relink style),
%     with each RQ tied to a §1 failure mode or framing axis.
%   - Ablation grid extended: now 5 rows including a fact-source
%     substrate ablation and a local-PPR ablation, each row tagged
%     with which §1 failure mode it targets.
%   - Main table and ablation numbers filled from current run logs.
%
% Data sources (carried from v1):
%   - Main graph baselines:
%     run_logs/baseline_qwen32b_graph_top200_gpt4omini_full_20260512
%     run_logs/hipporag_qwen32b_valid_graph_top200_full1000_20260513_r2
%   - NeocorRAG:
%     run_logs/neocorrag_aligned_k3_all32_full1000_20260514
%     run_logs/neocorrag_reader_gpt4omini_full1000_20260516
%   - HGRAG:
%     run_logs/hgrag_main3_32b_no_think_gpt4omini_full1000_20260517
%     run_logs/hgrag_nq_popqa_32b_no_think_gpt4omini_full1000_20260517
%   - RAPTOR:
%     run_logs/raptor_32b_no_think_gpt4omini_full1000_20260517
%   - EvidenceFlow:
%     run_logs/all32_etv3_etv4_pcec_gpt4omini_none_full1000_20260514
%   - Ablations:
%     run_logs/etv4_ablation_relation_coverage_all32_gpt4omini_full1000_20260516
%   - Dense-entry diagnostic:
%     run_logs/etv4_full1000_optimized_equivalence_20260511/reports/etv4_ablation6_dense_only_same_entry_full1000.json
%     run_logs/dense_entry_gpt4omini_reader_full1000_20260516
%
% =====================================================================

\section{Experiments}
\label{sec:experiments}

We design experiments around three research questions that mirror the
framing of \S\ref{sec:intro}. \textbf{RQ1.} How does \methodname{}
compare with no-retrieval, sparse, dense, hierarchical, OpenIE-based,
and path/chain-oriented baselines on multi-hop QA under a unified reader?
\textbf{RQ2.} Which components of \methodname{}---the fact-source
substrate, query-local PPR, relation-grounded expansion, and
coverage-aware reranking---contribute to retrieval and answer quality,
and do they target the two failure modes (bridge underweighting and
entity-neighborhood crowding) identified in \S\ref{sec:intro}?
\textbf{RQ3.} Does the evidence flow change which passages enter the
top-five reader prefix, beyond reordering dense neighbors? We report
paired bootstrap confidence intervals, operational statistics, and
query-level case studies in the Appendix.

\subsection{Experimental Setup}
\label{sec:experiments:setup}

\paragraph{Datasets.}
We evaluate on three widely used multi-hop question answering
benchmarks: HotpotQA \citep{yang2018hotpotqa}, 2WikiMultiHopQA
\citep{ho2020constructing}, and MuSiQue \citep{trivedi2022musique}.
The main comparison also reports NQ and PopQA as open-domain QA checks.
Following the evaluation protocol of HippoRAG~2
\citep{gutierrez2025hipporag}, we use a 1{,}000-query split per dataset
together with the corresponding retrieval corpus. This setting allows
us to report retrieval quality and downstream answer quality under the
same reader protocol.

\paragraph{Baselines.}
We group baselines into five categories that span the major paradigms
in retrieval-first multi-hop QA. \emph{(1) No Retrieval} reports the
reader's intrinsic capability without any retrieved context, using the
same GPT-4o-mini model we use downstream. \emph{(2) Sparse and Dense
Retrievers} retrieve passages by question--passage similarity without
any graph structure: BM25 \citep{robertson2009bm25} as a sparse
baseline, and the dense entry ranking from
NV-Embed-v2 \citep{lee2025nvembed} as the dense
baseline. \emph{(3) Hierarchical Tree Retrieval} uses RAPTOR
\citep{sarthi2024raptor}, which recursively clusters and summarizes
passages into a retrieval tree. \emph{(4) Graph and Hypergraph Retrievers} perform
non-iterative propagation over pre-built retrieval structures.
HippoRAG~2 \citep{gutierrez2025hipporag} extracts OpenIE facts and runs
personalized PageRank over a passage--phrase graph, while HGRAG
\citep{hgrag2024} constructs an entity hypergraph and applies
hypergraph diffusion over entity and passage granularities. \emph{(5) Path-- and
Chain-oriented Graph Retrievers} construct or mine reasoning chains at
query time: PropRAG \citep{wang2025proprag} performs beam search over a
proposition graph, and NeocorRAG \citep{peng2026neocorrag} mines
evidence chains from retrieved text and uses them to filter the reader
context. For methods that require language-model-assisted indexing,
chain extraction, or question-side analysis (HippoRAG~2, HGRAG,
PropRAG, NeocorRAG, and \methodname{}), we use Qwen3-32B with reasoning
traces disabled in the output. For final answer generation, all
methods use the same GPT-4o-mini reader.

\paragraph{Metrics.}
Following \citet{gutierrez2025hipporag}, we evaluate two subtasks.
Retrieval is measured by supporting-passage recall in the top-five
reader context, denoted R@5. End-to-end QA is measured by exact match
(EM) and token F1 against the reference answer. For ablations, we
additionally report All@5, the fraction of questions for which all
gold supporting passages appear in the top-five reader prefix. All
scores are reported as percentages.

\paragraph{Implementation details.}
All retrieval systems produce ranked passage lists, and the reader
consumes the same top-five reader context across methods. All methods
are evaluated on the same 1{,}000 questions and the same corpus
snapshot for each benchmark. HippoRAG~2 and PropRAG retrieve from graph
candidate pools capped at 200 passages.
\methodname{} builds a top-100 native candidate pool from the same
top-200 frontier, performs query-local propagation
(\S\ref{sec:method:retrieval}) over its fact-source substrate
(\S\ref{sec:method:substrate}), and passes the reranked list
(\S\ref{sec:method:enhancement}) to the same reader. Dense embeddings
are computed with NV-Embed-v2 for the fact-source substrate, and
question-side anchors $A(q)$ and requirements $B(q)$ are extracted by
Qwen3-32B with a fixed prompt. The reader uses temperature 0, and
EM/F1 are computed under the same answer normalization across methods.
For RAPTOR, we build recursive summary trees with Qwen3-32B summaries
and NV-Embed-v2 embeddings, retrieve tree nodes, and project selected
summary nodes back to descendant leaf passages before applying the same
top-five reader protocol.
Hyperparameters $\eta$ and $\lambda$ are tuned once on a small
development split and held fixed across datasets. Prompts and the full
hyperparameter table are provided in the Appendix.
Appendix~\ref{app:baseline-protocol} details how baseline outputs are
adapted to the shared top-five passage-reader protocol.

\subsection{Main Results (RQ1)}
\label{sec:experiments:main}

\begin{table*}[t]
\centering
\scriptsize
\setlength{\tabcolsep}{2pt}
\resizebox{\textwidth}{!}{%
\begin{tabular}{llccc ccc ccc ccc ccc c}
\toprule
& & \multicolumn{3}{c}{HotpotQA} &
\multicolumn{3}{c}{2WikiMultiHopQA} &
\multicolumn{3}{c}{MuSiQue} &
\multicolumn{3}{c}{NQ} &
\multicolumn{3}{c}{PopQA} &
MH Avg F1 \\
\cmidrule(lr){3-5}\cmidrule(lr){6-8}\cmidrule(lr){9-11}
\cmidrule(lr){12-14}\cmidrule(lr){15-17}
Category & Method & R@5 & EM & F1 & R@5 & EM & F1 & R@5 & EM & F1
& R@5 & EM & F1 & R@5 & EM & F1 &  \\
\midrule
\emph{No Retrieval}
 & GPT-4o-mini (no context)
 & -- & 23.9 & 31.4 & -- & 12.3 & 13.0 & -- & 7.5 & 11.7
 & -- & 29.2 & 37.8 & -- & 0.7 & 0.8 & 18.7 \\
\midrule
\emph{Sparse / Dense}
 & BM25
 & 72.9 & 49.4 & 61.0 & 65.8 & 45.1 & 49.0 & 47.8 & 20.5 & 28.0
 & 71.2 & 46.1 & 59.5 & 39.5 & 45.4 & 56.1 & 46.0 \\
 & NV-Embed-v2 (dense entry)
 & 93.4 & 59.5 & 73.0 & 75.9 & 55.0 & 60.5 & 68.3 & 35.2 & 46.2
 & 77.2 & 49.4 & 63.3 & 49.2 & 47.9 & 60.4 & 59.9 \\
\midrule
\emph{Hierarchical Tree}
 & RAPTOR
 & 90.9 & 59.0 & 72.0 & 71.7 & 49.6 & 54.4 & 65.2 & 31.5 & 41.5
 & 77.1 & 49.8 & 64.1 & 48.2 & 47.2 & 59.8 & 56.0 \\
\midrule
\emph{Graph / Hypergraph}
 & HippoRAG~2
 & 92.6 & 59.3 & 72.5 & 82.8 & 56.7 & 63.4 & 71.0 & 34.9 & 46.5
 & 79.8 & 50.3 & 63.4 & 49.6 & 47.7 & 60.3 & 60.8 \\
 & HGRAG
 & 94.5 & 60.6 & 73.3 & 76.7 & 54.5 & 60.2 & 69.5 & 36.8 & 48.1
 & 52.5 & 51.0 & 64.7 & 51.5 & 45.7 & 58.4 & 60.5 \\
\midrule
\emph{Path / Chain Graph}
 & PropRAG
 & 95.1 & 61.8 & 75.1 & 90.9 & 61.1 & 69.1 & \textbf{74.2} & 36.5 & 47.6
 & \textbf{81.1} & 50.1 & 63.7 & 56.5 & \textbf{49.4} & 61.6 & 63.9 \\
 & NeocorRAG\rlap{$^{\ast}$}
 & 95.0 & 62.1 & 74.8 & 88.9 & 59.8 & 66.5 & 73.1 & 37.1 & 48.1
 & -- & -- & -- & -- & -- & -- & 63.1 \\
\midrule
\emph{Ours}
 & \methodname{} & \textbf{96.5} & \textbf{62.8} & \textbf{75.6}
                 & \textbf{96.2} & \textbf{65.9} & \textbf{74.7}
                 & 73.5 & \textbf{37.2} & \textbf{48.9}
                 & 71.6 & \textbf{51.9} & \textbf{65.0}
                 & \textbf{65.0} & \textbf{49.4} & \textbf{62.2} & \textbf{66.4} \\
\bottomrule
\end{tabular}
\vspace{1pt}
}
\caption{Main results on three multi-hop QA benchmarks under a unified
GPT-4o-mini reader, with NQ and PopQA included as open-domain QA checks.
Methods are grouped by category. R@5 reports recall in the top-five
reader context, and EM/F1 report end-to-end answer quality from the
reader given that context. MH Avg F1 reports the multi-hop average over
HotpotQA, 2WikiMultiHopQA, and MuSiQue. $^{\ast}$NeocorRAG's published
evaluation does not include NQ or PopQA, and we omit these cells since
its query-time chain-mining prompt is not directly transferable to
open-domain settings.}
\label{tab:main-results}
\end{table*}

Table~\ref{tab:main-results} presents the main comparison. We discuss
each category in turn.

\paragraph{Versus no-retrieval and sparse--dense baselines.}
The no-retrieval row exposes how much of each dataset can be answered
from the reader's parameters alone, while BM25 and NV-Embed-v2
represent retrieval-without-structure. Across all three benchmarks,
\methodname{} substantially outperforms these baselines in both R@5
and F1. The clearest gap appears on 2WikiMultiHopQA, where moving from
the dense entry ranking to \methodname{} improves R@5 from 75.9 to
96.2 and F1 from 60.5 to 74.7. This is the dataset where multi-hop
bridges are most common, and the gap supports the
	\S\ref{sec:intro} claim that pointwise dense retrieval is built for
the opposite regime.

\paragraph{Versus hierarchical tree retrieval.}
RAPTOR performs competitively on HotpotQA and the open-domain checks,
but its multi-hop average F1 remains below the dense entry baseline and
the graph/path baselines. The largest gap again appears on
2WikiMultiHopQA, where its R@5 is 71.7 and F1 is 54.4. This suggests
that recursive summarization trees provide useful semantic abstraction
but do not by themselves enforce the bridge coverage needed by the
top-five reader prefix.

\paragraph{Versus graph and hypergraph retrievers.}
HippoRAG~2 shares our OpenIE indexing line but operates over a
passage--phrase graph, while HGRAG represents the hypergraph-based
GraphRAG route with entity nodes and passage hyperedges. \methodname{}
improves R@5 over HippoRAG~2 on all three multi-hop datasets, with the
largest gain again on 2WikiMultiHopQA (82.8 $\to$ 96.2). Compared with
HGRAG, \methodname{} gives higher F1 on HotpotQA and 2WikiMultiHopQA
and a comparable MuSiQue F1, while also improving NQ F1 under the same
reader. These comparisons position our fact-source propagation design
against both passage--phrase PPR and cross-granularity hypergraph
diffusion, although they are system-level comparisons rather than
substrate-only ablations.

\paragraph{Versus path-- and chain-oriented graph retrievers.}
PropRAG and NeocorRAG are the strongest baselines and outperform the
sparse/dense category, as expected for systems that explicitly
construct multi-hop structure. \methodname{} still gives the best R@5,
EM, and F1 on HotpotQA and 2WikiMultiHopQA. On MuSiQue, PropRAG holds
the highest R@5 while \methodname{} obtains the best EM and F1.
NeocorRAG saves reader contexts of average size 3.3--3.6 passages,
giving it a different precision--coverage trade-off; under the same
top-five reader, \methodname{} improves average F1 from 63.1 to 66.4
and average R@5 from 85.7 to 88.7. The result supports the central
claim of \S\ref{sec:intro}: a document-grounded propagation followed
by evidence-aware optimization can produce a competitive ranked
evidence list without query-time chain mining or path search.

\paragraph{Open-domain checks.}
NQ and PopQA are included as diagnostic open-domain checks rather than
as the primary target setting. Their trends are mixed in retrieval:
\methodname{} improves QA F1 on both datasets, but its NQ R@5 is lower
than the strongest dense and graph baselines. This is consistent with a
method designed primarily for multi-hop bridge recovery rather than
single-hop entity lookup.

\subsection{Ablation Study (RQ2)}
\label{sec:experiments:ablation}

\begin{table*}[t]
\centering
\small
\setlength{\tabcolsep}{4pt}
\begin{tabular}{lccc ccc ccc}
\toprule
& \multicolumn{3}{c}{HotpotQA} &
\multicolumn{3}{c}{2WikiMultiHopQA} &
\multicolumn{3}{c}{MuSiQue} \\
\cmidrule(lr){2-4}\cmidrule(lr){5-7}\cmidrule(lr){8-10}
Variant & R@5 & F1 & All@5 & R@5 & F1 & All@5 & R@5 & F1 & All@5 \\
\midrule
Full \methodname{} & \textbf{96.5} & \textbf{75.6} & \textbf{93.3}
& \textbf{96.2} & \textbf{74.7} & \textbf{89.5}
& \textbf{73.5} & 48.9 & \textbf{45.3} \\
\midrule
w/o local PPR \rlap{$^{\dagger}$}            & 93.4 & 73.0 & 87.1 & 75.9 & 60.5 & 49.3 & 68.3 & 46.2 & 37.2 \\
w/o relation expansion                        & 95.2 & 74.8 & 90.9 & 92.8 & 71.1 & 81.1 & 73.0 & \textbf{49.0} & 43.9 \\
w/o coverage reranking                        & 95.0 & 74.8 & 90.5 & 93.5 & 72.8 & 82.6 & 71.8 & 48.3 & 42.2 \\
\bottomrule
\end{tabular}
\caption{Component ablations under a unified GPT-4o-mini reader and a
single fixed configuration across datasets. All@5 measures whether the
top-five reader context contains \emph{all} gold supporting passages.
$^{\dagger}$Drops query-local PPR; the pipeline falls back to the dense
entry ranking. The HippoRAG~2 comparison in
Table~\ref{tab:main-results} provides a coarse proxy for moving from a
passage--phrase graph to the fact-source substrate, but it is not a
substrate-only ablation because propagation locality and the downstream
evidence-aware operators also differ.}
\label{tab:ablation}
\end{table*}

Table~\ref{tab:ablation} isolates the contribution of each component
inside \methodname{}. We tie each row to a claim from
\S\ref{sec:intro}. We do not claim a separate substrate-only ablation:
the HippoRAG~2 row in Table~\ref{tab:main-results} is a coarse
passage--phrase substrate proxy, while the full gap to \methodname{}
also includes query-local propagation and the evidence-aware operators
that run on the fact-source substrate.

\paragraph{Local PPR is the main retrieval engine.}
Removing query-local PPR (w/o local PPR) is equivalent to using the
dense entry ranking and produces the largest drop on 2WikiMultiHopQA,
where F1 decreases by 14.2 points and All@5 by 40.2 points. This is
consistent with the evidence flow $\pi_q$ being the core retrieval
signal, and with the two evidence-aware steps in
\S\ref{sec:method:enhancement} refining its ranking rather than creating it.

\paragraph{Relation-grounded expansion targets bridge underweighting.}
Removing relation-grounded expansion hurts 2WikiMultiHopQA most clearly
(R@5 $-$3.4, F1 $-$3.6, All@5 $-$8.4), consistent with its role of
admitting bridge and tail-entity passages that lie far from the query
anchors and receive little $\pi_q$-mass. This is consistent with the
bridge-underweighting failure mode identified in \S\ref{sec:intro}. The effect is smaller on
HotpotQA, where bridges are shorter, and mixed on MuSiQue: the F1
difference is within 0.1 point, while the full model still improves
R@5 and All@5. We therefore report MuSiQue as a retrieval-side gain
rather than a QA gain from expansion alone.

\paragraph{Coverage-aware reranking targets entity-neighborhood crowding.}
Removing coverage-aware reranking produces a more uniform retrieval
drop. R@5 decreases on all three datasets, and All@5 decreases by 2.8,
6.9, and 3.1 points on HotpotQA, 2WikiMultiHopQA, and MuSiQue
respectively. This is consistent with the entity-neighborhood
crowding failure mode identified in \S\ref{sec:intro}: graph
propagation alone can leave the head of the
ranked list concentrated around a single entity neighborhood, and the
noisy-OR coverage objective improves the final reader prefix by
preferring candidates that contribute to question-side requirements
not yet represented.

\subsection{Diagnostic Analysis (RQ3)}
\label{sec:experiments:analysis}

\begin{table*}[t]
\centering
\small
\setlength{\tabcolsep}{4pt}
\begin{tabular}{lccc ccc c}
\toprule
Dataset & \multicolumn{3}{c}{R@5} &
\multicolumn{3}{c}{All@5} &
Support gain \\
\cmidrule(lr){2-4}\cmidrule(lr){5-7}
& Dense entry & \methodname{} & $\Delta$ &
Dense entry & \methodname{} & $\Delta$ & \\
\midrule
HotpotQA & 93.4 & 96.5 & +3.1 & 87.1 & 93.3 & +6.2 & 7.2 \\
2WikiMultiHopQA & 75.9 & 96.2 & +20.4 & 49.3 & 89.5 & +40.2 & 45.1 \\
MuSiQue & 68.3 & 73.5 & +5.3 & 37.2 & 45.3 & +8.1 & 19.9 \\
\bottomrule
\end{tabular}
\caption{Retrieval-side diagnostic comparing the dense entry ranking
with \methodname{} under the same top-five reader prefix.
\emph{Support gain} reports the percentage of queries whose number of
retrieved gold supporting passages strictly increases in the top-five
list.}
\label{tab:flow-diagnostic}
\end{table*}

Table~\ref{tab:flow-diagnostic} addresses RQ3 by directly contrasting
the dense entry ranking with \methodname{} on the retrieval side.
\paragraph{Evidence flow changes the top-five evidence set.}
The largest gain appears on 2WikiMultiHopQA, where R@5 improves by
20.4 points and All@5 by 40.2 points. The support-gain column shows
that for 45.1\% of 2Wiki queries, \methodname{} retrieves strictly
more gold supporting passages in the top five than the dense entry
ranking. This is consistent with the core propagation step not simply
reordering the dense neighbors of the question, but changing
membership of the top-five list. The effect is smaller on HotpotQA,
where the dense baseline already retrieves most supporting passages,
and intermediate on MuSiQue.

\paragraph{Retrieval--QA coupling is not linear.}
Comparing Tables~\ref{tab:main-results} and \ref{tab:flow-diagnostic}
shows that retrieval and answer metrics are related but not identical.
On MuSiQue, PropRAG obtains the highest R@5 while \methodname{} obtains
the highest F1, indicating that better individual passage recall does
not always translate into better answer generation when questions
require longer or more fragmented evidence combinations. This is
consistent with the \S\ref{sec:intro} claim that a useful multi-hop
retriever should place a \emph{connected and covering} evidence set
into the reader prefix, not only recover individual gold passages.

\paragraph{Uncertainty and operational checks.}
Appendix~\ref{app:bootstrap-ci} reports query-paired bootstrap
intervals for the main deltas. The largest gains are stable under this
resampling. On 2WikiMultiHopQA, \methodname{} improves F1 over the dense
entry ranking by 14.2 points with a 95\% interval of [11.8, 16.6], and
improves R@5 by 20.4 points with an interval of [18.9, 21.9]. The
component-level intervals also separate the strong 2Wiki gains from the
smaller MuSiQue enhancement deltas, which we treat as retrieval-side
evidence rather than a robust QA gain. Appendix~\ref{app:operational-stats}
reports the corresponding graph and runtime statistics. Across the
three multi-hop datasets, the query-local graphs average about 200
passages and 1.4K--1.6K local edges, while the evidence-aware readout
adds 0.10--0.17 seconds per query excluding the reader. The offline
OpenIE substrate contains 68K--143K facts for the evaluated corpus
snapshots. Appendix~\ref{app:case-study} gives two query-level cases in
which the final readout admits a lower-ranked bridge passage into the
top-five prefix.
